\documentclass[12pt]{article}
\usepackage[utf8]{inputenc}
\usepackage[margin=1in]{geometry}
\usepackage{hyperref}
\usepackage{enumitem}
\usepackage{titlesec}
\usepackage{xcolor}
\usepackage{listings}

\titleformat{\section}{\large\bfseries\color{blue}}{}{0em}{}
\titleformat{\subsection}{\normalsize\bfseries\color{darkgray}}{}{0em}{}

\title{\textbf{Project Comprehensive Technical Guide: Satya \\ Generative AI Fake News \& Deepfake Detection}}
\author{Project Documentation Team}
\date{May 2026}

\begin{document}

\maketitle

\tableofcontents
\newpage

\section{Project Abstract}
The "Satya" project is an advanced, multi-modal verification platform designed to combat the rising tide of AI-generated misinformation and deepfakes. By integrating state-of-the-art Large Language Models (LLMs), Computer Vision, and a robust microservices architecture, the system provides users with real-time, evidence-based analysis of text, images, and videos. This document serves as a comprehensive technical manual for the project, detailing the architectural decisions, implementation strategies, and the underlying logic that powers the detection engine.

\section{The Problem Domain: Misinformation in the Age of GenAI}
In the modern digital landscape, the cost of generating high-quality fake content has dropped significantly due to the advent of Generative AI. 
\begin{itemize}
    \item \textbf{Textual Hallucinations:} AI models can generate convincing but entirely false news articles.
    \item \textbf{Deepfake Media:} Sophisticated video and audio manipulation can replicate public figures with high fidelity.
    \item \textbf{Information Overload:} Fact-checkers cannot keep up with the volume of content generated daily.
\end{itemize}
Our project addresses these issues by providing an automated, scalable tool that does not just "label" content but "explains" the reasoning behind its verdict.

\section{System Architecture: The Microservices Paradigm}
The project adopts a \textbf{decoupled microservices architecture}. This was a critical design choice to ensure that computationally expensive tasks (like video analysis) do not degrade the performance of lightweight tasks (like text analysis).

\subsection{1. API Gateway (The Orchestrator)}
The Gateway, built with \textbf{FastAPI}, serves as the brain's "pre-frontal cortex." It handles all incoming traffic, validates requests, and routes them to the appropriate worker services.
\begin{itemize}
    \item \textbf{Responsibility:} Security, CORS, rate limiting, and request routing.
    \item \textbf{Why FastAPI?} We chose FastAPI over traditional frameworks like Flask or Django because of its native support for \textbf{asynchronous programming} (ASGI), which is vital for handling concurrent file uploads and long-running I/O operations.
\end{itemize}

\subsection{2. Content Service (The Analytical Engine)}
This service handles text and image inputs. It contains the logic for:
\begin{itemize}
    \item \textbf{Web Scraping:} Extracting content from URLs provided by users.
    \item \textbf{OCR Pipeline:} Converting image text into machine-readable format.
    \item \textbf{GenAI Integration:} Orchestrating calls to high-reasoning LLMs.
\end{itemize}

\subsection{3. Video Service (The Background Specialist)}
Video analysis is done in an \textbf{asynchronous background environment}. 
\begin{itemize}
    \item \textbf{Mechanism:} When a video is uploaded, the Gateway stores it and pushes a job notification to a \textbf{RabbitMQ} queue.
    \item \textbf{State Storage:} The current status of the job (PENDING, PROCESSING, COMPLETED) is tracked in \textbf{Redis}.
    \item \textbf{Worker Logic:} A separate worker process consumes the queue, performs frame-by-frame analysis, and writes the final JSON report back to Redis.
\end{itemize}

\section{Technical Implementation Deep Dive}

\subsection{Text Verification Logic}
The verification isn't just a simple search. It follows a multi-step cognitive pipeline:
1. \textbf{De-contextualization:} Breaking down complex paragraphs into atomic factual claims.
2. \textbf{Query Generation:} Formulating optimized search queries to find evidence for or against each claim.
3. \textbf{Evidence Synthesis:} Comparing the claim against retrieved search results from diverse, credible sources.
4. \textbf{Conflict Resolution:} If sources disagree, the model evaluates the reputation of the sources and the recency of the information.

\subsection{Image \& OCR Processing}
Images are processed using \textbf{Tesseract OCR}. Since raw OCR output is often "noisy" (missing spaces, misread characters), we feed the OCR output into a secondary "Cleaning Model" (an LLM) that reconstructs the intended message before analysis.

\subsection{Video Deepfake Detection}
Our video service looks for \textbf{temporal and spatial anomalies}. This includes:
\begin{itemize}
    \item \textbf{Lip-Sync Discrepancies:} Checking if the audio waveform aligns with the mouth movements.
    \item \textbf{Edge Artifacts:} Detecting subtle blurring or flickering around the face area, which is common in GAN-generated content.
    \item \textbf{Metadata Analysis:} Checking for inconsistencies in file headers and compression patterns.
\end{itemize}

\section{Infrastructure \& DevOps}
\begin{itemize}
    \item \textbf{Docker:} Every service is containerized, ensuring that the development environment matches the production environment perfectly.
    \item \textbf{RabbitMQ:} Acts as the message broker, ensuring that even if the worker service crashes, the jobs are not lost (Persistence).
    \item \textbf{Redis:} Provides sub-millisecond latency for job tracking, which is essential for a smooth frontend polling experience.
\end{itemize}

\section{The Future of Satya}
Our roadmap includes:
1. \textbf{Blockchain Logging:} Storing verification reports on a ledger to prevent tampering with historical truth.
2. \textbf{Browser Extension:} Real-time analysis as users browse social media feeds.
3. \textbf{Audio-only Analysis:} Extending detection to AI-generated voice clones and "Vishing" attacks.

\section{Conclusion}
The Satya project demonstrates that fighting AI with AI is the only viable path forward. By combining a modern web stack with robust infrastructure and cutting-edge machine learning models, we have built a tool that empowers users to distinguish truth from fiction in the digital age.

\end{document}
